\section{BC: Base-Field-Normalized Split-Pencil Census}
\label{sec:gf-work-bc}

\paragraph{Status.}
This section separates the part of BC that is already proved in the raw
CAP25 package from the part that remains a genuine census input.  The slope
elimination, interpolation-lattice, determinantal-representation, automatic
divisibility, boundary-profile, and base-field-floor reductions below are
PROVED relative to \cite{CAP25v13}.  The final interior split-pencil count is
still CONJECTURAL; its finite deployed form requires exact row constants and
cannot be replaced by an unspecified polynomial loss.

\begin{definition}[BC parameters and split locators]
\label{def:gf-bc-parameters}
Let \(C=\RS[\F,D,K]\), where \(D\subseteq\B\subseteq\F\), \(|D|=n\),
\(|\B|=b\), \(|\F|=q\), and
\[
        \Lambda_D(X)=\prod_{x\in D}(X-x).
\]
Fix an agreement \(m\ge K\), and put
\[
        \omega=n-m,\qquad w=m-K .
\]
An interior balanced profile is an integer \(d_1\) satisfying
\[
        w+2\le d_1\le \left\lfloor {n-K+1\over2}\right\rfloor,
        \qquad d_2=n-K+1-d_1 .
\]
The excluded endpoint \(d_1=w+1\) is the quotient/prefix boundary profile and
belongs to Q, not to BC.

A monic polynomial \(W\) is \(D\)-split of degree \(\omega\) if
\(W\mid\Lambda_D\) and \(\deg W=\omega\).  Thus \(W=\Lambda_{D\setminus T}\)
for a unique \(m\)-subset \(T\subseteq D\).
\end{definition}

\begin{definition}[determinantal datum and BC pencil count]
\label{def:gf-bc-pencil-count}
For a word \(U:D\to\F\), let
\[
        M_U=\{(W,N)\in\F[X]^2:\ W(x)U(x)=N(x)\ \text{for every }x\in D\}
\]
and set
\[
        \operatorname{wdeg}(W,N)=
        \max\{\deg W,\deg N-(K-1)\}.
\]
A determinantal datum of profile \((d_1,d_2)\) is a quadruple
\[
        \mathfrak g=(W_1,N_1,W_2,N_2)
\]
arising from a shifted weak-Popov basis \(g_i=(W_i,N_i)\) of some \(M_U\), with
\(\operatorname{wdeg}(g_i)=d_i\).  Equivalently, after choosing a unit
\(\gamma\in\F^\times\),
\[
        W_1N_2-W_2N_1=\gamma\Lambda_D,\qquad \gcd(W_1,W_2)=1,
        \qquad d_1+d_2=n-K+1.
\]
For an interior profile define the split-pencil count
\[
\begin{aligned}
        \mathrm{BC}_{\mathfrak g}(m,d_1)
        :=\#\Bigl\{(A,B):\;&
        \deg A\le \omega-d_1,\quad
        \deg B\le \omega-d_2,\\
        &W_{A,B}:=AW_1+BW_2
        \text{ is monic, }D\text{-split, and }\deg W_{A,B}=\omega
        \Bigr\}.
\end{aligned}
\]
The monic condition fixes the scalar normalization of the pencil.
\end{definition}

\begin{definition}[base-field floor at an interior profile]
\label{def:gf-bc-base-floor}
For \(d_1\) as in \cref{def:gf-bc-parameters}, put
\[
        m_{d_1}:=K-1+d_1 .
\]
The exact base-field floor forced by the identity-prefix construction at this
profile is
\[
        \underline M_{\B}(d_1)
        :=
        \binom{m_{d_1}}m
        \left\lceil \binom n{m_{d_1}} b^{-(d_1-1)}\right\rceil .
\]
Its density-model version, used for asymptotic statements, is
\[
        M_{\B}(d_1)
        :=
        \binom{m_{d_1}}m
        \binom n{m_{d_1}} b^{-(d_1-1)} .
\]
\end{definition}

\begin{proposition}[slope elimination]
\label{prop:gf-bc-slope-elimination}
Let \((u,v)\) be an affine received line, and let \(T\subseteq D\) have
\(|T|=m\).  For \(0\le r<w\) define
\[
        S_r(h,T)=\sum_{x\in T}{h(x)x^r\over\Lambda_T'(x)},\qquad
        \alpha(T)=(S_r(u,T))_{r<w},\quad
        \beta(T)=(S_r(v,T))_{r<w}.
\]
Then \((u+zv)|_T\) extends to a polynomial of degree \(<K\) if and only if
\[
        \alpha(T)+z\beta(T)=0 .
\]
If \(T\) is not common, meaning \((\alpha(T),\beta(T))\ne(0,0)\), then it
supports at most one finite slope.  Consequently, after common-support and
tangent cells have been paid, the number of unpaid bad finite slopes at
agreement \(m\) is bounded by the number of non-common rank-one supports,
namely supports with \(\alpha(T)\) and \(\beta(T)\) proportional and
\(\beta(T)\ne0\).
\end{proposition}

\begin{proof}
The Lagrange interpolant of \(h:T\to\F\) has degree \(<K\) if and only if its
top \(w=m-K\) coefficients vanish.  The displayed functionals are a
unitriangular change of coordinates from those top coefficients, so this is
equivalent to \(S_r(h,T)=0\) for every \(r<w\).  Applying this to \(h=u+zv\)
and using linearity gives the vector equation.  If \(\beta(T)=0\), either
\(\alpha(T)=0\) and every slope is carried by the common-support cell, or no
slope is carried by \(T\).  If \(\beta(T)\ne0\), any solution \(z\) is forced
by one nonzero coordinate of \(\beta(T)\), and it exists exactly when
\(\alpha(T)\) and \(\beta(T)\) are proportional.  Choosing one non-common
support for each unpaid bad slope gives the injection.
\end{proof}

\begin{proposition}[lattice and split-pencil reduction]
\label{prop:gf-bc-lattice-split}
For a word \(U:D\to\F\), the agreement-support census is
\[
\begin{aligned}
        \mathrm{cen}(U;m)
        =\#\Bigl\{(W,N)\in M_U:\;&
        W\text{ monic},\ \deg W=\omega,\
        \operatorname{wdeg}(W,N)\le\omega,\\
        &W\mid\Lambda_D,\quad W\mid N
        \Bigr\}.
\end{aligned}
\]
If the minimal shifted degree of \(M_U\) is an interior profile \(d_1\), and
\((g_1,g_2)\) is a shifted weak-Popov basis of profile \((d_1,d_2)\), then
every census element has the form
\[
        (W,N)=A g_1+B g_2,\qquad
        \deg A\le\omega-d_1,\quad \deg B\le\omega-d_2 .
\]
Moreover, for any such element with \(W\) \(D\)-split, the divisibility
\(W\mid N\) is automatic.  Hence the balanced interior census of \(U\) is
bounded by, and with the monic exact-degree normalization equals, the
split-pencil count of \cref{def:gf-bc-pencil-count}.
\end{proposition}

\begin{proof}
If \(T\) is a valid agreement support and \(c\in\RS[\F,T,K]\) is the unique
degree-\(<K\) codeword agreeing with \(U\) on \(T\), then
\(W=\Lambda_{D\setminus T}\) and \(N=Wc\) satisfy the displayed conditions.
Conversely, if \((W,N)\) satisfies them, write \(N=Wc\); the shifted-degree
cap gives \(\deg c<K\), and \(W(U-c)=0\) on \(D\), so \(U\) and \(c\) agree
on the \(m\)-point support \(D\setminus Z(W)\).

The module \(M_U\) is free of rank two with basis \((1,\widehat U)\) and
\((0,\Lambda_D)\).  A shifted weak-Popov basis has predictable degrees and row
degree sum \(n-K+1\), so an element of shifted degree at most \(\omega\) has
the displayed representation with the displayed degree bounds.  Finally, if
\(W\mid\Lambda_D\), then \(W\) is squarefree and all its roots lie in \(D\).
At each root \(x\) of \(W\), the defining relation of \(M_U\) gives
\(N(x)=W(x)U(x)=0\), so \(W\mid N\).
\end{proof}

\begin{proposition}[boundary profile is Q]
\label{prop:gf-bc-boundary-is-q}
The profile \(d_1=w+1\) of the split-pencil formulation is exactly the
quotient/prefix divisor problem: for a prefix value \(z\), the census is the
number of monic degree-\(\omega\) divisors of \(\Lambda_D\) with the prescribed
top \(w+1\) coefficients.  Therefore BC starts at \(d_1=w+2\); counting
\(d_1=w+1\) inside BC would double-count Q and would use the wrong scale.
\end{proposition}

\begin{proof}
This is the Q-unification theorem of the raw package.  For the identity-prefix
witness \(U_z\), the pair \((1,P_z)\), where
\[
        P_z=X^m+\sum_{j\le w}z_jX^{m-j},
\]
has shifted degree \(w+1\).  The companion row is obtained from the division
\(\Lambda_D=P_zQ_z+R_z\).  In the balanced branch, predictable degrees force
every census element to be \(Q_z+A\) after monic normalization, with
\(\deg A\le\omega-w-1\), so the split condition is precisely
\((Q_z+A)\mid\Lambda_D\).  This is the prefix-fiber problem and is assigned to
Q.
\end{proof}

\begin{proposition}[base-field floor is unavoidable]
\label{prop:gf-bc-base-floor}
Assume \(m_{d_1}+d_1\le n\).  There exists a word whose shifted lattice has
minimal profile \(d_1\) and whose size-\(m\) balanced census is at least
\[
        \underline M_{\B}(d_1)
        =
        \binom{m_{d_1}}m
        \left\lceil \binom n{m_{d_1}} b^{-(d_1-1)}\right\rceil .
\]
Consequently no correct BC model may be normalized only at the challenge-field
scale \(\binom n\omega q^{1-w}\) when \(q\) is larger than \(b\).
\end{proposition}

\begin{proof}
Use the level-\(m_{d_1}\) identity-prefix witness with prefix depth
\(d_1-1\):
\[
        P_{z'}^{(m_{d_1})}
        =
        X^{m_{d_1}}+\sum_{j\le d_1-1} z'_j X^{m_{d_1}-j},
        \qquad z'\in\B^{d_1-1}.
\]
The same shifted-degree computation as in the boundary profile gives minimal
degree \(d_1\) under the hypothesis \(m_{d_1}+d_1\le n\).  By pigeonhole, the
heaviest prefix value has at least
\[
        \left\lceil \binom n{m_{d_1}} b^{-(d_1-1)}\right\rceil
\]
level-\(m_{d_1}\) supports.  Each such support gives a distinct codeword
agreeing with the witness on \(m_{d_1}\) points, and each such codeword
contributes \(\binom{m_{d_1}}m\) size-\(m\) sub-supports to the binomial
agreement moment.  Distinct level supports give distinct codewords, so the
contributions do not collide.
\end{proof}

\begin{conjecture}[BC, exact census form]
\label{conj:gf-bc-exact-census}
For every smooth multiplicative or circle row, every agreement \(m\) in the
normalized band, and every primitive interior determinantal datum
\(\mathfrak g\) after tangent, common-support, common-GCD, quotient-pullback,
subfield-confined, extension-line, boundary-Q, and bounded-chart cells have
been paid, one has
\[
        \mathrm{BC}_{\mathfrak g}(m,d_1)
        \le
        R_{BC}(n)\,
        \max\left(
        1,\ M_{\B}(d_1),\ \binom n\omega q^{1-w}
        \right)
\]
uniformly over the interior range
\[
        w+2\le d_1\le \left\lfloor {n-K+1\over2}\right\rfloor .
\]
The asymptotic target is \(R_{BC}(n)=e^{o(n)}\), or any explicitly bounded
polynomial loss used only with logarithmic agreement reserve.  The finite
adjacent target replaces \(R_{BC}(n)\) by row-wise constants
\(R_{BC,\mathrm{row}}(d_1)\) and by exact integer ceilings for the three model
terms; the resulting sum over profiles must fit inside the printed one-step
margins.
\end{conjecture}

\begin{remark}[finite packet required for BC]
\label{rem:gf-bc-finite-packet}
A finite BC packet for one deployed row should print, for each interior
profile \(d_1\):
\[
\begin{array}{ll}
\text{row data} & (\F,\B,D,K,n,m,\omega,w,q,b),\\
\text{profile} & (d_1,d_2),\\
\text{paid removals} &
\text{tangent/common, GCD, quotient, subfield, extension, Q, bounded charts},\\
\text{datum} & (W_1,N_1,W_2,N_2),\
W_1N_2-W_2N_1=\gamma\Lambda_D,\\
\text{pencil equations} & \deg A\le\omega-d_1,\ \deg B\le\omega-d_2,\
AW_1+BW_2\mid\Lambda_D,\\
\text{normalization} & AW_1+BW_2\text{ monic of degree }\omega,\\
\text{count} & \mathrm{BC}_{\mathfrak g}(m,d_1)
\le R_{BC,\mathrm{row}}(d_1)\max(1,\underline M_{\B}(d_1),
\lceil\binom n\omega q^{1-w}\rceil),\\
\text{ledger use} & \text{first-match insertion into }U(a_0+1),\\
\text{status} & \text{PROVED / CONDITIONAL / CONJECTURAL / COUNTEREXAMPLE}.
\end{array}
\]
Without these row constants, BC can support only the reserve/asymptotic closure,
not the deployed adjacent theorem.
\end{remark}

\begin{remark}[proof routes and citations]
\label{rem:gf-bc-routes}
The reductions above leave a single interior census problem: split divisors of
\(\Lambda_D\) lying in a two-generator affine pencil.  The plausible routes are:
\begin{enumerate}[label=\textup{(\alph*)}]
\item \textbf{Master flatness.}  Prove a base-field split-locator flatness
theorem for affine conditions on monic split divisors, with quotient and
common-divisor components removed.  This would imply the asymptotic BC bound
with reserve.  The formulation should be aligned with the master flatness
target in \cite{CAP25v13}.
\item \textbf{Determinantal incidence and monodromy.}  Treat
\(AW_1+BW_2\mid\Lambda_D\) as a point-count problem on the determinantal
incidence variety attached to
\[
        W_1N_2-W_2N_1=\gamma\Lambda_D .
\]
The expected obstruction strata are exactly the base-field floor, quotient
pullbacks, common factors, and bounded charts; a primitive excess beyond these
would be a new obstruction floor.
\item \textbf{Exchange-compression.}  Try to show that exchanging roots toward
graded or quotient-stabilized configurations cannot decrease the number of
split points in a codimension condition.  A proof would turn the known floor
strata into extremizers; a counterexample would falsify the current frontier
model by producing a primitive split-pencil family.
\item \textbf{Rung transfer for dyadic rows.}  Use the tower
\(x\mapsto x^2\) to separate periodic mass exactly and prove a bounded-loss
transfer for primitive residuals.  A constant loss per rung is only polynomial
asymptotically, but the finite deployed packet must still track the bit cost
at every rung.
\item \textbf{Singular-bucket cleanup.}  For overdetermined or low-deficiency
charts, use Kronecker pencil structure and Berlekamp--Massey/Pad'e
two-generator normal forms to show that identically split charts collapse to
tangent or already-paid cells.  These tools should be cited explicitly if this
route is promoted from strategy to theorem.
\end{enumerate}
\end{remark}

\begin{remark}[overclaiming guardrails]
\label{rem:gf-bc-overclaiming}
The following statements are not proved by the reductions in this section:
\begin{enumerate}[label=\textup{(\roman*)}]
\item the bound in \cref{conj:gf-bc-exact-census} for primitive interior
split pencils;
\item any finite adjacent safe certificate at \(a_0+1\);
\item replacement of the base-field floor by a challenge-field floor;
\item merging BC with Q or SP in a point estimate; and
\item using \(e^{o(n)}\) or polynomial losses in the finite deployed rows
without printed constants fitting the row margins.
\end{enumerate}
These items should remain labelled CONJECTURAL or CONDITIONAL until a replayable
certificate or a proof with explicit constants is supplied.
\end{remark}
